\documentclass[12pt]{article}

\usepackage[T2A]{fontenc}
\usepackage[utf8]{inputenc}
\usepackage[russian]{babel}
\usepackage{geometry}
\usepackage{amsmath}
\usepackage{hyperref}

\geometry{a4paper, margin=2cm}

\title{Текст к презентации по курсовой работе}
\author{И.\,А. Буканов}
\date{2026}

\begin{document}
\maketitle

\section*{Общий план на 15 минут}

Доклад должен быть численно конкретным и заранее закрывать ожидаемые вопросы. Главная линия рассказа:
я построил пайплайн измерения SBF по JWST/F150W, откалибровал абсолютную SBF-величину по независимой TRGB-шкале, проверил устойчивость через leave-one-out и сравнил результат с опубликованной HST/F160W SBF-калибровкой Jensen et al. (2015).

Ориентировочный тайминг:
\begin{itemize}
    \item 1.5 минуты --- постановка задачи и данные;
    \item 4 минуты --- физика SBF и технический пайплайн: модель, остатки, $P(k)$;
    \item 2 минуты --- ошибки и калибровка $\bar{M}_{150}$;
    \item 5 минут --- основные графики, дельты и заранее закрытые спорные точки;
    \item 1.5 минуты --- ограничения и выводы.
\end{itemize}

\section*{Слайд 1. Тема}

Сказать: работа посвящена измерению расстояний до галактик ближней Вселенной методом флуктуаций поверхностной яркости. Конкретно здесь я работаю с данными JWST/NIRCam в фильтре F150W и пытаюсь получить практически применимую калибровку SBF в этом фильтре.

\section*{Слайд 2. Цель}

Сказать: цель не просто получить одну величину для одной галактики, а построить полный воспроизводимый пайплайн. На входе --- изображения JWST F150W/F090W, на выходе --- $\bar{m}_{150}$, ошибка, цветовая калибровка $\bar{M}_{150}$ и расстояние.

Ключевая формулировка: независимая опорная шкала здесь --- TRGB. Она нужна не потому, что мы хотим заново переписать TRGB-расстояния, а потому что SBF требует абсолютной калибровки. После калибровки SBF можно применять там, где TRGB измерять сложнее.

\section*{Слайд 3. Метод SBF}

Сказать: в неразрешённой галактике пиксели содержат разное число ярких звёзд. Поэтому даже после вычитания гладкой модели остаются статистические флуктуации. Амплитуда этих флуктуаций уменьшается с расстоянием.

Главная формула:
\begin{equation*}
    \mu = \bar{m}_{150} - \bar{M}_{150}(F090W-F150W).
\end{equation*}
Здесь $\bar{m}_{150}$ измеряется из изображения, а $\bar{M}_{150}$ нужно откалибровать.

\section*{Слайд 4. Данные}

Сказать: использованы публичные i2d-продукты JWST/NIRCam программы GO-3055 из MAST. F150W используется для SBF, F090W нужен для цвета. Всего обработано 14 галактик. Основная калибровочная выборка состоит из 12 объектов: два объекта имеют реальные предупреждения моделирования и не входят в clean-fit. Предупреждение о размере файла не считается научным флагом, потому что кадры были проверены визуально.

Литературные модули расстояния взяты из TRGB-SBF Project IV. Для внешнего сравнения используются опубликованные HST/F160W SBF-измерения Jensen et al. (2015).

Здесь же сразу закрыть вопрос про цефеиды и сверхновые: это не игнорирование хороших методов, а ограничение конкретной выборки. Большинство объектов --- ранние галактики, для которых нет однородной прямой Cepheid/SN Ia шкалы. Если смешать разные внешние источники, мы внесём дополнительную систематику. Поэтому TRGB используется как единая опорная шкала, а Jensen/F160W --- как внешняя SBF-проверка.

\section*{Слайд 5. Пайплайн}

Сказать: пайплайн состоит из фона, маски компактных источников, изофотной модели галактики, остаточного изображения, измерения спектра мощности и подгонки вида
\begin{equation*}
    P(k)=P_0E(k)+P_1.
\end{equation*}
Здесь $E(k)$ задаётся PSF и маской, $P_0$ --- амплитуда флуктуаций, $P_1$ --- белый шум. Из $P_0$ вычитается вклад неразрешённых источников $P_r$, после чего получается SBF-сигнал.

Нужно подчеркнуть: я не просто считал число, а проверял остатки глазами, потому что для SBF качество гладкой модели критично.

Фраза для перехода: теперь подробно, что именно значит ``гладкая модель'', потому что это один из главных источников возможной систематики.

\section*{Слайд 6. Как строилась модель галактики}

Сказать: первичная модель Серсика использовалась технически --- чтобы заполнить NaN и центральную область, где изофоты могут быть нестабильны. Финальная модель для SBF не является просто Серсиком. Основная гладкая модель строится по эллиптическим изофотам: центр, эллиптичность и позиционный угол могут меняться с радиусом.

Ключевой момент: изофоты не ограничивались искусственно, а шли настолько далеко, насколько сходились устойчиво. Это нужно, чтобы не получить видимый шов между внутренней изофотной частью и внешней достроенной частью. Такой шов сразу проявился бы как крупномасштабный остаток и испортил бы низкие частоты в $P(k)$.

Сказать прямо: модель проверялась визуально по трём изображениям --- исходный кадр, модель и остаток image--model. Это важно, потому что формально корректный fit может быть непригоден для SBF, если в остатке остаётся крупномасштабная структура.

\section*{Слайд 7. Как получался остаток для SBF}

Сказать: после вычитания модели остаётся остаточное изображение. Но SBF нельзя мерить по нему напрямую, потому что там есть шаровые скопления, фоновые галактики и звёзды. Поэтому применялись две маски: первичная маска компактных источников до построения модели и дополнительная маска положительных компактных остатков после вычитания модели.

Важно сформулировать: эта маска не вырезает сам SBF-сигнал. SBF --- это распределённая статистическая мощность, размазанная PSF и присутствующая во всей области галактики. Маска вырезает отдельные компактные объекты, которые иначе дали бы лишнюю мощность и завысили бы $P_0$.

Области измерения фиксированы: внутреннее и внешнее кольца $8.2$--$16.4''$ и $16.4$--$32.8''$. Это делает сравнение между галактиками более однородным.

\section*{Слайд 8. Как фиттился спектр мощности $P(k)$}

Сказать: SBF измеряется не как средняя яркость остатка, а в частотном пространстве. Для каждой области считается спектр мощности $P(k)$. Модель спектра:
\begin{equation*}
    P(k)=P_0E(k)+P_1.
\end{equation*}

Здесь $E(k)$ --- форма ожидаемого SBF-сигнала. Она определяется PSF и маской области: физически SBF-флуктуации свёрнуты с PSF, а маска меняет спектр. $P_0$ --- амплитуда флуктуаций, именно она нужна для $\bar{m}_{150}$. $P_1$ --- почти белый шум, который не несёт SBF-информации.

Рабочее окно $k=0.01$--$0.25$ cycles/pixel выбрано не произвольно. Очень низкие частоты чувствительны к недовычтенной модели галактики. Очень высокие частоты чувствительны к пиксельному шуму и деталям PSF. Поэтому фит идёт в диапазоне, где форма $P(k)$ согласована с $E(k)$ и не доминирует ни модельный фон, ни чистый шум.

После fit применяется поправка:
\begin{equation*}
    P_{\rm fluc}=P_0-P_r,
\end{equation*}
где $P_r$ --- вклад неразрешённых шаровых скоплений и фоновых галактик. Только $P_{\rm fluc}$ переводится в $\bar{m}_{150}$.

\section*{Слайд 9. Ошибки}

Сказать: ошибка измеренной $\bar{m}_{150}$ считается как квадратичная сумма независимых вкладов:
\begin{equation*}
    \sigma_{\bar{m}}^2 =
    \sigma_{P(k)}^2 + \sigma_{\rm bkg}^2 +
    \sigma_{\rm PSF}^2 + \sigma_{P_r}^2 + \sigma_{\rm ext}^2.
\end{equation*}

Главный вклад обычно даёт $\sigma_{P(k)}$, то есть неопределённость подгонки спектра мощности. PSF меньше, но методологически важен, потому что входит в $E(k)$. $P_r$ нужен, чтобы не принять слабые шаровые скопления и фоновые галактики за настоящую SBF-флуктуацию.

Отдельно сказать: ошибка расстояния больше ошибки $\bar{m}_{150}$, потому что к ней добавляется ширина калибровки $\sigma_{\rm int}\simeq0.049$ mag.

\section*{Слайд 10. Калибровка $\bar{M}_{150}$}

Сказать: для каждой галактики я беру
\begin{equation*}
    \bar{M}_{150}=\bar{m}_{150}-\mu_{\rm TRGB}.
\end{equation*}
Затем строю зависимость от цвета:
\begin{equation*}
    \bar{M}_{150}=(-3.454\pm0.013)+(1.17\pm0.21)[(F090W-F150W)-0.40].
\end{equation*}
Число 0.40 mag не имеет отдельного физического смысла: это центрирование около середины диапазона цветов, чтобы нуль-пункт был устойчивым.

Сразу проговорить, почему точки не лежат идеально на прямой: это реальные галактики, у них есть ошибки измерения, ошибки TRGB, ошибки цвета, различия звёздных популяций и остаточный внутренний разброс $\sigma_{\rm int}$. Серая полоса на графике показывает именно внутренний разброс калибровки, а не ошибку среднего значения. Поэтому наличие точек вне центральной линии не является провалом модели.

\section*{Слайд 11. Зачем цвет}

Сказать: без цвета можно было бы считать $\bar{M}_{150}$ постоянной. Но это хуже. По leave-one-out постоянная модель даёт разброс 0.102 mag, а линейная цветовая калибровка --- 0.073 mag. В расстояниях это примерно улучшение с 4.7\% до 3.4\%, то есть примерно 28\% относительно постоянной калибровки.

Квадратичная модель внутри выборки выглядит почти так же, как линейная, но на leave-one-out ведёт себя хуже: 0.084 mag против 0.073 mag. Поэтому на таком числе объектов линейная модель наиболее честная.

\section*{Слайд 12. Сравнение с TRGB}

Сказать: после калибровки я восстанавливаю модуль расстояния
\begin{equation*}
    \mu_{150}=\bar{m}_{150}-\bar{M}_{150}(C).
\end{equation*}
Сравнение с TRGB даёт
\begin{equation*}
    \langle\mu_{150}-\mu_{\rm TRGB}\rangle=0.003\pm0.019~\mathrm{mag},
\end{equation*}
WRMS равен 0.061 mag. В расстояниях это средний сдвиг около 0.15\% и разброс около 2.8\%.

Сразу пояснить график: чёрная линия --- это равенство SBF и TRGB, а не линия подгонки. Поэтому вопрос ``почему не все точки лежат на прямой'' неправильно трактует график. Отклонения ожидаемы из-за ошибок отдельных измерений, ошибок TRGB и внутреннего разброса калибровки. Важно, что среднее смещение практически нулевое.

Отдельно про NGC 1380: она действительно визуально заметна, потому что лежит ниже линии равенства. Но численно это $\Delta\mu=-0.082\pm0.067$ mag, то есть около $1.2\sigma$. Это не статистически значимое расхождение и не основание исключать объект. Такая точка показывает реальный разброс метода.

\section*{Слайд 13. Leave-one-out}

Сказать: это более строгая проверка, потому что каждая галактика по очереди исключается из калибровки, а потом её расстояние предсказывается по остальным объектам. Разброс получается около 0.073 mag, то есть примерно 3.4\% по расстоянию.

Фраза для защиты результата: это не полностью независимая космологическая шкала, но это проверка того, что результат не держится на одной галактике и что калибровка применима к объектам вне обучающей точки.

\section*{Слайд 14. Сравнение с Jensen et al. 2015}

Сказать: я сравниваю свою JWST/F150W SBF-шкалу с опубликованной HST/F160W SBF-шкалой Jensen et al. Для пяти общих галактик среднее смещение:
\begin{equation*}
    \langle\mu_{150}-\mu_{160,\rm Jensen}\rangle=-0.063\pm0.067~\mathrm{mag}.
\end{equation*}
Это значит, что мои расстояния в среднем немного меньше, примерно на 2.9\%.

Сразу закрыть вопрос о расхождении: это не выглядит противоречием. Во-первых, сравниваются разные фильтры и инструменты: JWST/F150W против HST/F160W. Во-вторых, Jensen 2015 связан со старой Cepheid/optical-SBF шкалой, а моя калибровка завязана на TRGB. Jensen 2025 обсуждает, что TRGB-основанная шкала для Virgo/Fornax оказывается примерно на 3--4\% короче старой Cepheid/optical-SBF шкалы. Поэтому найденные $-2.9\%$ имеют правильный знак и порядок величины.

\section*{Слайд 15. Диагностика систематик}

Сказать: я отдельно проверил остатки против цвета, расстояния, разности внутренней и внешней областей, ошибки $P(k)$, PSF и вклада $P_r/P_0$. На 12 объектах нельзя делать сильные выводы о корреляциях: статистика слишком мала. Но можно проверить отсутствие очевидного провала пайплайна. Явного одного параметра, который объясняет весь остаточный разброс, не видно.

Здесь же закрыть вопрос про NGC 1399: пайплайн по ней формально прошёл, но это не значит, что точка обязана быть точной. У NGC 1399 ошибка $\bar{m}_{150}$ равна 0.098 mag, то есть существенно больше, чем у хороших объектов. Это центральная массивная галактика с более сложной крупномасштабной структурой и заметной разностью между внутренней и внешней областями. Поэтому она получает малый вес и не определяет калибровку.

\section*{Слайд 16. Итоги}

Сказать коротко:
\begin{itemize}
    \item измерены $\bar{m}_{150}$ для 14 галактик JWST GO-3055;
    \item построена F150W SBF-калибровка по TRGB-шкале;
    \item цвет уменьшает разброс расстояний с 0.102 до 0.073 mag;
    \item относительно TRGB среднее смещение равно $0.003\pm0.019$ mag;
    \item относительно Jensen/F160W найдено $-0.063\pm0.067$ mag, что согласуется с известным различием шкал.
\end{itemize}

\section*{Короткие формулировки для вопросов после доклада}

Эти ответы теперь уже проговорены в основном рассказе. Если вопрос всё равно зададут, отвечать коротко и численно.

\subsection*{Почему на графике расстояний не все точки лежат на прямой?}

Потому что линия --- это равенство с TRGB, а не результат подгонки. Отклонения ожидаемы: есть ошибка измерения $\bar{m}_{150}$, ошибка TRGB, внутренняя ширина калибровки $\sigma_{\rm int}$ и реальные различия звёздных популяций. Численно разброс 0.061 mag соответствует 2.8\% по расстоянию, а среднее смещение $0.003\pm0.019$ mag статистически совместимо с нулём.

\subsection*{Почему NGC 1380 заметно ушла на графике?}

Для NGC 1380 получается $\mu_{150}-\mu_{\rm TRGB}=-0.082$ mag при ошибке 0.067 mag. Это около $1.2\sigma$, то есть не статистически значимая катастрофа. В сравнении с Jensen/F160W она тоже смещена, но опять в пределах объединённых ошибок. Объект не исключён, потому что визуально остатки и формальные ошибки пайплайна нормальные; это вклад в реальный разброс калибровки.

\subsection*{Почему у NGC 1399 большие дельты, если пайплайн прошёл?}

NGC 1399 формально прошла, но это не значит, что она обязана иметь малую ошибку. У неё $\bar{m}_{150}=28.125\pm0.098$ mag, то есть ошибка почти в три раза больше, чем у лучших объектов. Также у неё большая разность внутренних и внешних областей. Это массивная центральная галактика с более сложной крупномасштабной структурой, поэтому результат получает малый вес. При этом относительно Jensen/F160W для неё нет большого конфликта: разность почти нулевая.

\subsection*{Почему нет калибровки по цефеидам или сверхновым Ia?}

Потому что для этой конкретной выборки нет однородной прямой Cepheid/SN Ia шкалы. Большинство объектов --- ранние галактики, где цефеиды не являются естественным методом. Сверхновые Ia есть не для всех и не дают однородную таблицу для этих же галактик. Если смешать разные внешние источники, мы добавим систематику. Поэтому для этой работы честнее использовать TRGB как основную независимую шкалу и Jensen/F160W как внешнюю SBF-проверку.

\subsection*{Почему сравнение с Jensen/F160W не совпадает идеально?}

Во-первых, это другой телескоп и другой фильтр: JWST/F150W против HST/F160W. Во-вторых, Jensen 2015 опирается на старую SBF/Cepheid-связанную шкалу, а здесь калибровка завязана на TRGB. В Jensen 2025 как раз обсуждается, что TRGB-основанная шкала для Virgo/Fornax короче старой Cepheid/optical-SBF шкалы на 3--4\%. Наше среднее смещение $-2.9\%$ имеет правильный знак и порядок величины.

\subsection*{Почему калибровка через литературные расстояния не круговая?}

Это стандартная процедура калибровки индикатора расстояний. TRGB-расстояния используются не как конечный результат, а чтобы определить абсолютную величину SBF. После этого калибровку можно применять к другим галактикам, где TRGB получить труднее. Для проверки устойчивости внутри этой же выборки используется leave-one-out.

\subsection*{Что именно улучшено на 28\%?}

Не абсолютная точность всей внегалактической шкалы расстояний, а качество восстановления расстояний внутри этой SBF-калибровки. Если считать $\bar{M}_{150}$ постоянной, leave-one-out разброс равен 0.102 mag, что соответствует примерно 4.7\% по расстоянию. Если учитывать цвет $F090W-F150W$, разброс падает до 0.073 mag, то есть примерно до 3.4\%. Относительное уменьшение ошибки составляет около 28\%.

\subsection*{Почему основной fit строится по 12, если обработано 14 галактик?}

Обработано 14 галактик. Две из них имеют реальные предупреждения устойчивости моделирования: NGC 4621 и NGC 4697. Поэтому основная чистая калибровка строится по 12 объектам. Предупреждение о размере файла не считается научным флагом, потому что соответствующие кадры были проверены визуально.

\end{document}
